\medskip
	On considère la suite $(u_{n})$ définie par $u_{0}=5$ et, pour tout entier naturel $n$ :
	$$
	u_{n+1}=2+\ln \left(u_{n}^{2}-3\right).
	$$

	On admet que cette suite est bien définie.

	\subsection*{Partie A : Exploitation de programmes Python}

	\begin{enumerate}
		\item Recopier et compléter le script Python ci-dessous pour que \texttt{suite(k)} qui prend en paramètre un entier naturel \texttt{k}, renvoie la liste des $k$ premières valeurs de la suite $(u_{n})$.

		\medskip

		\textbf{Remarque :} On précise que, pour tout réel strictement positif \texttt{a}, \texttt{log(a)} renvoie la valeur du logarithme népérien de \texttt{a}.















		\begin{center}
		\begin{ttfamily}
		\begin{tabularx}{0.6\linewidth}{|X|}\hline
			def suite(k):\\
			\quad	L = []\\
			\quad	u = 5\\
			\quad	for i in range(......):\\
			\qquad	L.append(u)\\
			\qquad	u=............\\
			\quad	return(......)\\ \hline
		\end{tabularx}
		\end{ttfamily}
		\end{center}
		
		\item On a exécuté \texttt{suite(9)} ci-dessous. Émettre deux conjectures : l'une sur le sens de variation de la suite $(u_{n})$ et l'autre sur son éventuelle convergence.

\begin{center}
\begin{ttfamily}
		
\begin{tabularx}{\linewidth}{|X|}\hline
>{}>{}>{} suite(9) \\
$\left[\right.$ 5, 5.091042453358316, 5.131953749864703,\\
5.150037910978289, 5.157974010229213, 5.1614456706362954,\\
5.162962248594583, 5.163624356938671, 5.163913344065642$\left.\right]$\\ \hline

\end{tabularx}
\end{ttfamily}
\end{center}

		\item On a ensuite créé la fonction \texttt{mystere(n)} donnée ci-dessous et exécuté \linebreak \texttt{mystere(10000)}, ce qui a renvoyé \texttt{1}.

		Cet affichage contredit-il la conjecture émise sur le sens de variation de la suite $(u_{n})$ ? Justifier.

\begin{center}
		\begin{ttfamily}

\begin{tabularx}{\linewidth}{|X|}\hline
			def mystere(n):\\
			\quad L = suite(n)\\
			\quad	c = 1\\
			\quad	for i in range(n - 1):\\
			\qquad	if L[i] $>$ L[i + 1]:\\
			\qquad \quad c = 0\\
			\quad return c\\ \hline

\end{tabularx}
		\end{ttfamily}
		\end{center}

		\begin{center}
		\begin{ttfamily}

\begin{tabularx}{\linewidth}{|X|}\hline
		>{}>{}>{} mystere(10000)\\
		1\\ \hline
\end{tabularx}

		\end{ttfamily}
		\end{center}
	\end{enumerate}

\subsection*{Partie  B : Étude de la convergence de la suite ( $\boldsymbol{u}_{\boldsymbol{n}}$ )}

On considère la fonction $g$ définie sur $[2 ~;~+\infty [$ par :

\[g(x)=2+\ln \left(x^{2}-3\right)\]

	On admet que $g$ est dérivable sur $\left[2 ~;~+\infty\left[\right.\right.$ et on note $g'$ sa fonction dérivée.

\begin{enumerate}
	\item Démontrer que la fonction $g$ est croissante sur $[2 ~;~+\infty[$.

	\item 
		\begin{enumerate}
		\item Démontrer par récurrence que, pour tout entier naturel $n$ :

\[4 \leqslant u_{n} \leqslant u_{n+1} \leqslant 6.\]

		\item En déduire que la suite $\left(u_{n}\right)$ converge.
		\end{enumerate}
	\end{enumerate}

\subsection*{Partie C : Étude de la valeur de la limite}

On considère la fonction $f$ définie sur [ 2 ; $+\infty$ [ par :

\[f(x)=2+\ln \left(x^{2}-3\right)-x.\]

On admet que $f$ est dérivable sur $[2 ~;~+\infty[$ et on note $f'$ sa fonction dérivée.

On donne le tableau de variations de $f$ suivant. On ne demande aucune justification.










\[\begin{tablvar}[6em]{2}
\hline
x & 2 && 3 && +\infty\\
\hline
\variations{\mil{f(x)} & \bas{0} && \haut{\ln(6)-1} &&  \bas{-\infty}} 
\hline
\end{tablvar}\]

	\begin{enumerate}
		\item \begin{enumerate}
			\item Montrer que l'équation $f(x)=0$ admet exactement deux solutions sur \linebreak $[2 ~;~+\infty[$ que l'on notera $\alpha$ et $\beta$ avec $\alpha<\beta$.

			\item Donner la valeur exacte de $\alpha$ et une valeur approchée à $10^{-3}$ près de $\beta$.
		\end{enumerate}

		\item On note $\ell$ la limite de la suite $(u_{n})$.

		Justifier que $f(\ell)=0$ et déterminer $\ell$.
	\end{enumerate}
